% Auto-generated from Chapter-04-Haralds-Training-Oath.md
% Source: C:\Users\PCGAME\Desktop\story&universes\chain-weapon-story\finished-manuscript\manuscriptbaseforchapters\Chapter-04-Haralds-Training-Oath.md

\chapter{Harald's Training Oath}
\renewcommand{\CurrentChapterNum}{4}
\POV{\glyphAisen}{Aisen}

"Why do you want to learn?" Harald asked.

Aisen was ten years old now, and they were standing in the clearing behind the tavern, a space his father had carved out specifically for training. It was early morning, the sun not yet high enough to burn off the mist that clung to the grass. Harald stood with his spear resting against his shoulder, not threatening but present. He asked this question every few days now, variations on the same theme.

"Because I want to be strong," Aisen said, which was the answer he'd given last week.

"No," Harald said immediately. "Try again."

"Because I want to be able to fight people."

"No."

"Because I need to be able to defend myself."

"Better," Harald conceded. "But still not the real answer. What do you \emph{actually} want?"

Aisen stood in the morning cold and thought about this. He'd been training with Harald for two years now, ever since the robbery. Two years of being taught to move, to position himself, to understand the space around him. Two years of his body learning things his mind sometimes struggled to articulate. But this question—this question was different. This question insisted on truth rather than acceptable answers.

"I don't know," Aisen said finally, honestly.

Harald smiled. It was a rare expression on his face, used sparingly and therefore meaningful. "That's the first true thing you've said. Don't you want to know?"

"Yes," Aisen said.

"Then we start there."

They walked, not toward the training ground but toward the road. Harald moved with the economy of someone who'd learned that wasted motion was wasted life. Aisen, who'd been trained to notice such things, found himself mirroring the movement without quite meaning to—shoulders back, but not stiffly; stride measured; attention distributed around them rather than focused forward.

They walked in silence until they reached the bridge that crossed the river running south of town. The water was high with early summer runoff, moving fast, carrying debris from the mountains.

"Do you see the river?" Harald asked.

"Yes."

"Why does it flow the way it does? Why does it go around the rock instead of through it?"

Aisen looked at the river. At the place where it split around a large stone that had fallen perhaps years ago. "Because it has to," Aisen said.

"Yes. But why?"

Aisen thought about this. "Because the rock is hard and the water is soft. So the water has to move."

"And is the water strong or weak?"

"Weak," Aisen said.

"Then why are we walking on a bridge and not swimming through a canyon?"

Aisen understood then. "Because the water is patient. Because it keeps moving and eventually it wins."

Harald nodded. "Now. If you were the river, and I was the rock, what would be your advantage?"

"I don't know."

"Yet. You will." Harald turned away from the bridge and began walking back toward the tavern. "Your father employs the same principle in the tavern. He doesn't fight against customer resistance. He flows around it. He finds where they're soft and moves there. He eventually gets what he wants because his patience is stronger than their stubbornness."

"Is that what I should want to learn?" Aisen asked. "Patience?"

"That's \emph{part} of what you should want to learn," Harald said. "But not all. Come."

They returned to the training ground.

Over the course of the day, Harald taught him something new. Not a new technique—new understanding. He showed Aisen how a man's shoulders telegraphed his intent. How the way a person breathed changed when they were committed to an action versus when they were merely considering it. How the space around a person—the positioning of their feet, the bend of their knees—revealed whether they had options or whether they were already committed to the path they'd chosen.

"Every movement is a conversation with the world around you," Harald said, standing in front of Aisen with a staff. "Your body is speaking. The environment is speaking. With practice, you learn to listen."

He demonstrated, moving slowly, showing how his shoulders moved before his arms moved, how his feet repositioned before his entire body committed. It was like watching someone think out loud. It was like watching his father make decisions in the tavern, except made manifest through the body instead of through words.

"Now you try," Harald said.

Aisen tried. His body was clumsy, unused to moving with such deliberation. But Harald didn't criticize. He simply watched and made observations: "You're thinking about being strong. Stop that. Think about being efficient. There's difference."

"What's the difference?"

"Strong breaks things. Efficient moves things. We want efficiency. We want to spend less energy than we're forcing others to spend. We want to understand our advantage before we commit to using it."

Through the afternoon, through exercises that seemed initially pointless and gradually began to reveal their architecture, Aisen learned. His hands blistered from the staff work. His muscles ached. But more than that, his mind began to understand something about decision-making that he hadn't grasped before: the moment between stimulus and response was where all choice happened.

When evening came, and Harald was satisfied with the day's work, they sat on the edge of the training ground, and Harald did something unusual. He removed a small length of leather from his wrist, something Aisen had seen him wear but never questioned.

"I've had this since I was younger than you," Harald said. "Since I was trained by someone who understood that knowledge arranged correctly becomes power."

He handed it to Aisen.

It was braided leather, worn soft from years of handling. The braiding was complex—three strands woven in a pattern that seemed almost mathematical in its precision.

"He's gone now," Harald continued. "And I wear it to remember that what he taught me mattered. That learning mattered. That the choice to pay attention mattered." He took the leather back and re-secured it to his wrist. "One day, when you've learned enough, I'll give you a knot of your own. But not yet. You're still asking why."

"Is asking why a bad thing?" Aisen asked.

"No," Harald said. "It's the only thing. Strength without asking why is just violence. Efficiency without asking why is just cruelty. What separates people who matter from people who don't is that the people who matter ask why before they decide."

Aisen thought about his father. About how his father's questions always preceded his actions. About how his father never gave orders that were orders—he always asked questions that led people to understand why the order was necessary.

"Is that what this is?" Aisen asked. "Is this about making me a person who matters?"

Harald considered this. "This is about giving you the tools to choose who you become. Whether you choose to matter—that's your decision."

Three months later, they formalized the apprenticeship. It wasn't a ceremony, exactly. More of a mutual understanding that his training would no longer be casual but structured. That Harald would teach him not just the movements of combat but the philosophy underneath the movements. That Aisen would be his student not just in the afternoons but in everything—the way Harald moved through the tavern, the way he read people, the way he made choices.

It happened on a night when his father came to watch them train.

Harald and Aisen were working through a pattern—a series of movements that built one upon the other, each movement creating an opportunity for the next, each opportunity building toward something larger. Aisen was learning to execute them with less thought now. His body was beginning to understand.

When they finished, his father approached them.

"He's quick," Aisen's father said to Harald.

"He thinks," Harald replied. "That matters more than quick."

"I know," Aisen's father said. He turned to Aisen. "Do you want this? Do you want to be Harald's student?"

Aisen understood this was the formal question. The one where his answer would matter. "Yes," he said.

"Then understand what you're choosing," his father said. "Learning comes before natural talent. The work comes before the reward. And there's no guarantee that the reward will justify the work. But if you commit anyway, commit thoroughly."

"I commit," Aisen said, and he meant it.

His father turned to Harald. "Then he's yours to teach. But remember—he's still a boy. There are things he shouldn't bear yet."

"I know," Harald said.

Three days later, Harald gave him a gift. A spear—not a toy spear, but a real one, worn and scarred, the wood smooth from years of use. The blade wasn't particularly sharp, but the balance was perfect.

"This is mine," Harald said. "I'm giving it to you because you've given me your commitment. Keep it. Learn it. Treat it like it's part of you. When you understand the weight of it, when you understand how it moves and what it wants to do, you'll understand a lot of things that will help you navigate the world."

Aisen held the spear and felt its weight settle into his hands. It was heavier than he expected but also more balanced than expected. The wood was warm to the touch, warmed by years of Harald's hands.

"Will I ever have to use it?" Aisen asked.

"I don't know," Harald said. "I hope not. But if you do, you'll use it because you understand it, not because you're afraid. And that's the difference between a tool and a weapon."

That night, Aisen lay in his bed and turned the knowledge over in his mind. He was a student now. Not just a child learning casual skills but an apprentice to someone who knew more than he did. It felt like a responsibility and also like a gift—like someone had seen something in him that mattered, had seen potential worth investing in.

He thought about rivers and rocks. About the moment between stimulus and response. About the difference between strong and efficient.

He thought about the fact that his father had asked him if he was choosing this, had made clear that nothing was being forced. That the commitment was his own.

And he thought about Harald's hands, moving with such economy and precision, speaking a language that Aisen was only beginning to understand. The language of choices made well, of intentions clear, of understanding the world precisely enough to move through it without creating unnecessary harm.

He still didn't know exactly why he wanted to learn. But he was beginning to understand.

He wanted to understand why people did what they did. He wanted that knowledge to matter. He wanted to be someone whose presence changed things, the way his father's presence changed things in the tavern—quietly, efficiently, making space for others to become their better selves.

That was worth learning for.

That was worth the blisters and the aching muscles and the endless questions. That was worth committing to.